% Generated draft: requires mathematical and source review.
\documentclass[11pt]{article}
\usepackage[margin=25mm]{geometry}
\usepackage{fontspec}
\setmainfont{Latin Modern Roman}
\usepackage{amsmath,amssymb,mathtools,unicode-math}
\setmathfont{Latin Modern Math}
\usepackage{enumitem,needspace,xurl,hyperref}
\setlength{\parindent}{0pt}
\setlength{\parskip}{5pt}
\setlength{\emergencystretch}{3em}
\newcommand{\N}{\mathbb{N}}
\newcommand{\Z}{\mathbb{Z}}
\newcommand{\Q}{\mathbb{Q}}
\newcommand{\R}{\mathbb{R}}
\newcommand{\C}{\mathbb{C}}
\newcommand{\F}{\mathbb{F}}
\newcommand{\A}{\mathbb{A}}
\newcommand{\PP}{\mathbb{P}}
\newcommand{\E}{\mathbb{E}}
\newcommand{\eps}{\varepsilon}
\newcommand{\dd}{\,\mathrm{d}}
\begin{document}
\section*{UnsolvedMath ID: 30005907}
\textbf{Source problem code:} catalogue entry
\subsection{Definitions and mathematical statement}
Let \(p\) be a prime and let \(G\) be an elementary abelian \(p\)-group of rank \(d\), meaning that \(d\) is a nonnegative integer and
\[
G\cong (\Z/p\Z)^d.
\]
Let \(X\) be a finite CW complex equipped with a continuous action of \(G\). The action is free if, for every \(x\in X\) and \(g\in G\), the equality \(g\cdot x=x\) implies that \(g\) is the identity element. For a coefficient field \(k\), write \(H_n(X;k)\) for ordinary singular homology and define the total Betti number by
\[
b_k(X):=\sum_{n\in\Z}\dim_k H_n(X;k).
\]
This sum is finite because \(X\) is a finite CW complex.

The supplied formulation of the conjecture attributed to Gunnar Carlsson asks whether, for every prime \(p\), every rank \(d\), every such group \(G\), and every finite CW complex \(X\) described by the source as ``non-trivial,''
\[
G\text{ acts freely on }X
\quad\Longrightarrow\quad
b_k(X)\geq 2^d.
\]
The coefficient-field hypothesis and the meaning of ``non-trivial'' remain unresolved in the supplied record. In particular, the record does not specify the characteristic of \(k\) or its relationship to \(p\). The displayed implication therefore records the intended statement template, rather than asserting a conjecture quantified over all coefficient fields. No connectedness or cellularity hypothesis is supplied, and none is added here.
\subsection{Short English statement}
The question attributed to Gunnar Carlsson asks whether a free action of an elementary abelian \(p\)-group of rank \(d\) on a nontrivial finite CW complex forces its total homology dimension over \(k\) to be at least \(2^d\). The supplied record leaves the coefficient-field condition and the meaning of ``nontrivial'' unspecified, so these require clarification before the statement is mathematically definitive.
\subsection{Sources}
\begin{itemize}
\item Local UnsolvedMath registry: \texttt{data/unsolvedmath/problems.json}, permanent ID \(30005907\), title ``Carlsson's Homology Bound for Free Group Actions,'' exact source problem code \texttt{catalogue entry}; collection: Oberwolfach Reports Open Problems.
\item Supplied report citation: \emph{Hochschild (Co)Homology and Applications} (2024), \emph{Oberwolfach Reports}, pp.~1097--1162, DOI \texttt{10.4171/owr/2024/20}. Supplied link: \url{https://doi.org/10.4171/owr/2024/20}. The extracted text identifies Petter Andreas Bergh's contribution ``Complexes with small homology'' and attributes the conjecture to Gunnar Carlsson.
\item The extracted report text says that the conjecture was proposed in reference [3] and that Carlsson settled the case \(p=2\), \(d\leq 3\), in reference [4]. The bibliographic entries for [3] and [4] are not supplied. Its assertion that the conjecture remained open is qualified as being at the time of writing.
\item Supplied literature source: Bergh, ``homology over finite tensor categories,'' \url{https://doi.org/10.4171/JNCG/564}.
\item Supplied literature source: Ruping and Stephan, ``nonrealizable complexes,'' \url{https://doi.org/10.1016/j.jpaa.2022.107023}.
\item Literature assessment dated exactly \texttt{2026-08-22}: the supplied curation labels the problem ``partially solved'' and reports that Carlsson's topological bound is proved in several low-rank and special cases but remains open in general, while the associated algebraic conjecture is disproved for odd primes of rank at least eight. This is supplied dated curation derived from the OWR record, not an independently verified assessment. The associated algebraic conjecture is not formulated in the supplied material.
\end{itemize}
\end{document}
